\subsubsection{Disjoint Sparse Table}

\paragraph{Idea intuitiva.}
Estructura de datos estatica que permite responder consultas sobre cualquier rango $[l, r]$ en tiempo $O(1)$ tras un precomputo de $O(n \log n)$, pero a diferencia de la Sparse Table clasica, \textbf{no requiere idempotencia} (funciona para cualquier operacion asociativa: suma, multiplicacion, XOR, matrices, etc.). La idea: en cada nivel $k$, el arreglo se divide en bloques de tamano $2^{k+1}$. En cada bloque con punto medio \texttt{mid}, se precomputa un sufijo hacia la izquierda desde $\text{mid}-1$ y un prefijo hacia la derecha desde \texttt{mid}. Para cualquier consulta $[l, r]$, si se toma $k = \lfloor \log_2(l \oplus r) \rfloor$ (el bit mas significativo donde $l$ y $r$ difieren), $l$ y $r$ caen en el mismo bloque pero en lados opuestos de su punto medio. Asi, el rango $[l, r]$ se descompone exactamente en dos subrangos disjuntos: $[l, \text{mid}-1]$ y $[\text{mid}, r]$, sin solapamiento alguno.

\paragraph{Propiedades clave (invariantes).}
\begin{itemize}\setlength\itemsep{0pt}
	\item \texttt{st[k][i]}: en el nivel $k$, si $i < \text{mid}$ almacena la acumulacion en el sufijo $[i, \text{mid}-1]$; si $i \ge \text{mid}$, almacena la acumulacion en el prefijo $[\text{mid}, i]$.
	\item Identificacion de nivel en $O(1)$: la condicion de pertenencia al mismo bloque y lados opuestos del punto medio equivale a que $l$ y $r$ difieran en el bit $k$, por lo que $k = 31 - \text{\_\_builtin\_clz}(l \oplus r)$.
	\item La operacion debe ser \textbf{asociativa}; no requiere ser conmutativa ni idempotente (si no es conmutativa, se respeta el orden estricto: $\text{op}(\text{st}[k][l], \text{st}[k][r])$).
\end{itemize}

\paragraph{Codigo clave.}
Construccion y consulta $O(1)$ con Disjoint Sparse Table:
\begin{verbatim}
struct DisjointSparseTable {
    int n, LOG;
    vector<vector<int>> st;
    DisjointSparseTable(const vector<int>& a) {
        n = a.size();
        LOG = (n <= 1 ? 1 : 32 - __builtin_clz(n - 1) + 1);
        st.assign(LOG, a);
        for (int k = 1; k < LOG; ++k) {
            int half = 1 << k;
            for (int mid = half; mid < n; mid += 2 * half) {
                for (int i = mid - 2; i >= mid - half; --i)
                    st[k][i] = min(a[i], st[k][i + 1]);
                int lim = min(n, mid + half);
                for (int i = mid + 1; i < lim; ++i)
                    st[k][i] = min(st[k][i - 1], a[i]);
            }
        }
    }
    int query(int l, int r) {
        if (l == r) return st[0][l];
        int k = 31 - __builtin_clz(l ^ r);
        return min(st[k][l], st[k][r]);
    }
};
\end{verbatim}

\paragraph{Complejidad.}
\begin{itemize}\setlength\itemsep{0pt}
	\item Construccion: $O(n \log n)$ en tiempo.
	\item Consulta: $O(1)$ en el peor caso (una sola operacion \texttt{op} y calculo de \texttt{clz}).
	\item Memoria: $O(n \log n)$ enteros.
\end{itemize}

\paragraph{Donde tocar para modificar.}
\begin{itemize}\setlength\itemsep{0pt}
	\item \textbf{Cualquier operacion asociativa}: sustituir \texttt{min} por suma (\texttt{+}), multiplicacion modular, XOR (\texttt{\^{}}) o multiplicacion de matrices $O(M^3)$.
	\item \textbf{Operaciones no conmutativas}: asegurar el orden de los argumentos: $\text{op}(\text{st}[k][l], \text{st}[k][r])$.
	\item \textbf{Tabulacion de logaritmos}: para evitar \verb|__builtin_clz|, se puede precomputar un arreglo \texttt{lg} con $\lfloor \log_2 i \rfloor$.
\end{itemize}

\paragraph{Trampas y errores frecuentes.}
\begin{itemize}\setlength\itemsep{0pt}
	\item \textbf{Caso base $l = r$}: si $l = r$, entonces $l \oplus r = 0$, para el cual \verb|__builtin_clz(0)| es comportamiento indefinido en C++. Es indispensable tratar $l = r$ por separado retornando \texttt{st[0][l]}.
	\item \textbf{Indices 0-based}: la propiedad del bit mas significativo requiere que los indices sean 0-based ($0 \le l \le r < n$).
	\item \textbf{Limites en el ultimo bloque}: en el bucle derecho, verificar que \texttt{lim} no supere $n$ (\texttt{min(n, mid + half)}).
\end{itemize}

\paragraph{Explicacion linea por linea.}
\textbf{Estructura DisjointSparseTable:}
\begin{itemize}\setlength\itemsep{0pt}
	\item \verb|struct DisjointSparseTable { int n, LOG; vector<vector<int>> st; };| -- encapsula la estructura.
	\item \verb|LOG = (n <= 1 ? 1 : 32 - __builtin_clz(n - 1) + 1);| -- calcula $\lceil \log_2 n \rceil + 1$ niveles.
	\item \verb|st.assign(LOG, a);| -- inicializa cada nivel con los valores base del arreglo.
	\item \verb|for(int k=1; k<LOG; ++k)| -- itera sobre cada nivel con bloques de tamano $2^{k+1}$.
	\item \verb|int half = 1 << k;| -- la mitad del bloque actual; el punto de division es \texttt{mid}.
	\item \verb|for(int mid = half; mid < n; mid += 2 * half)| -- recorre los puntos medios de cada bloque.
	\item \verb|st[k][i] = min(a[i], st[k][i + 1]);| -- precalcula acumulaciones hacia la izquierda (sufijos).
	\item \verb|st[k][i] = min(st[k][i - 1], a[i]);| -- precalcula acumulaciones hacia la derecha (prefijos).
	\item \verb|int query(int l, int r)| -- resuelve la consulta en tiempo constante $O(1)$.
	\item \verb|if(l == r) return st[0][l];| -- evita $l \oplus r = 0$ y retorna el elemento unico.
	\item \verb|int k = 31 - __builtin_clz(l ^ r);| -- nivel cuyo punto medio separa a $l$ de $r$.
	\item \verb|return min(st[k][l], st[k][r]);| -- combina el sufijo $[l, \text{mid}-1]$ y el prefijo $[\text{mid}, r]$.
\end{itemize}
\textbf{Programa principal:}
\begin{itemize}\setlength\itemsep{0pt}
	\item \verb|DisjointSparseTable dst({3, 1, 4, 1, 5, 9, 2, 6});| -- construye la DST para el arreglo.
	\item \verb|cout << dst.query(0, 4) << "\n";| -- consulta el rango $[0, 4]$, resultado es 1.
	\item \verb|cout << dst.query(2, 7) << "\n";| -- consulta el rango $[2, 7]$, resultado es 1.
\end{itemize}

\paragraph{Codigo.}
Ver \texttt{cpp.tex}, seccion \textit{Disjoint Sparse Table}.

\vspace{0.5em}

